\chapter{Background and Literature Review}
\label{chap:background}

This chapter positions the thesis against five parts of the work that each contribute to the overall problem, which is to deploy a robust GeoFM on board a resource-constrained satellite.
Each part is treated in turn, before a closing synthesis states what combination of them remains unaddressed.

\section{Geospatial Foundation Models}

GeoFMs extend the foundation-model paradigm from natural images to Earth observation, pre-training a single large backbone on diverse, multi-modal EO archives so that many downstream tasks can be solved by fine-tuning a small task-specific head rather than training a full model from scratch.
TerraMind\cite{jakubik2025terramind} is representative of this generation: a ViT-based, multimodal GeoFM pre-trained across nine geospatial modalities at scale, explicitly designed to produce transferable semantic representations rather than task-specific features.
Studies compressing and deploying such models, discussed further below, treat this generalization capacity as the primary reason to prefer a GeoFM over a bespoke architecture -- but that capacity has, to date, mostly been demonstrated on data drawn from the same sensor family the model was pre-trained on, leaving open how it behaves under a genuine sensor change.

\section{Domain Adaptation}

Domain adaptation addresses the general problem of a model trained on a source distribution being applied to a target distribution with different input statistics but unchanged semantics \cite{kouw2019introductiondomainadaptationtransfer}.
Within satellite on-board AI specifically, Du et al.
\ \cite{du2024domainadaptation} demonstrated domain adaptation for satellite-borne cloud detection by retraining task heads on real, though limited, labeled imagery from the Kanyini mission, with the encoder frozen.
This line of work was extended on-orbit by Du et al.
\ \cite{du2025onorbitdemonstrationgeospatialfoundation}, who deployed a compressed Prithvi-family GeoFM on Kanyini and the IMAGIN-e payload, again adapting via real-label head retraining rather than encoder-level invariance training, and measuring the domain gap only indirectly through downstream task performance on source and target benchmark datasets that differ in geography and content as well as sensor.
This is the central limitation this thesis's data design addresses: without co-registered source/target pairs of the same scene, a measured performance drop cannot distinguish sensor-induced covariate shift from confounding differences in what the two datasets actually depict.

\section{Physics-Based Sensor Degradation Simulation}

Because labeled data from a satellite sensor does not exist before the satellite is launched, and accumulates slowly afterward for rare events, physics-based simulation of a target sensor's characteristics from an existing labeled source (typically Sentinel-2) has become a standard tool for pre-mission model development.
The simulator developed for the ESA $\Phi$-lab OrbitalAI challenge \cite{longepe2024orbitalai} models this transformation as a sequence of physical operations in radiance space -- spectral and spatial resampling to the target sensor's band configuration and ground sampling distance, band-to-band misalignment consistent with push-broom acquisition geometry, signal-to-noise degradation, and optical blur modeled through the sensor's modulation transfer function -- calibrated against real flight data from related missions.
This simulator underlies both the degradation pipeline used in this thesis and, independently, the pre-training and benchmarking datasets used by $\Phi$satNet (Section~\ref{sec:phisatnet-background}), making it a shared, validated tool within this specific research context rather than a novel contribution of either work.

\section{Knowledge Distillation Across Architectures}

Knowledge distillation \cite{hinton2015distilling} compresses a large teacher model into a smaller student by training the student to reproduce the teacher's output distribution, or, in feature-based variants, its intermediate representations, rather than training the student only against hard labels.
Most applications of distillation to on-board EO models compress within the same architecture family or distill into a same-family, smaller model, and treat robustness to sensor conditions as a property that transfers passively if the teacher itself is robust.
This thesis instead distills across architecture families -- from a ViT-based teacher to a CNN student -- a necessary step given that ViT inference is not supported on the $\Phi$-sat-2 on-board hardware (Section~\ref{sec:onboard-ai}), and applies the domain-invariance objective directly to the student during distillation rather than relying on passive inheritance from an already-invariant teacher.

\section{On-board AI for Earth Observation}
\label{sec:onboard-ai}

On-board AI processing has moved from single-purpose CNNs, exemplified by $\Phi$-sat-1's cloud-detection application\cite{giuffrida2021varphi}, toward more general and more capable models, as surveyed by Barretta et al.
\ \cite{barretta2026onboardaireview}.
The hardware reality constraining this progress is specific and severe: the $\Phi$-sat-2 CogniSAT-XE1 payload's Intel Movidius Myriad 2 VPU \cite{myriad2} is legacy hardware whose supported operator set (through OpenVINO 2020.3, ONNX opset 8--11) lacks dynamic-axes support, making ViT inference structurally infeasible regardless of model size, and imposes hard memory (approximately 512\,MB) and model-size (250\,MB) limits under float16-only inference.
These constraints, documented in detail by the same-lab $\Phi$satNet effort, are the concrete justification for why this thesis treats cross-architecture distillation, rather than direct deployment or same-family compression, as a requirement rather than a design choice.

\subsection{$\Phi$satNet}
\label{sec:phisatnet-background}

$\Phi$satNet is a U-Net-derived, multi-task-pretrained GeoFM developed within the same research group specifically for on-board execution on $\Phi$-sat-2, using a purpose-built pre-training corpus and benchmark suite built on the same OrbitalAI simulator described above.
By training a compact, non-transformer architecture from scratch, $\Phi$satNet avoids the ViT-hardware-incompatibility problem entirely, and its authors report that it matches or exceeds unconstrained transformer-based GeoFMs on several on-board tasks.
This is a different design philosophy from the one pursued here, not a competing solution to the same problem: $\Phi$satNet optimizes for a hardware-native architecture trained on a curated, mission-specific task suite, while this thesis bets that adapting an existing, diversely pre-trained GeoFM generalizes to new tasks with fewer labels than training a mission-specific model affords.
Chapter~\ref{chap:results-discussion} tests this bet directly through a few-shot comparison against $\Phi$satNet.

\section{Synthesis}

None of the five bodies of work above, individually or in combination, isolates sensor-induced covariate shift using co-registered source/target data, trains encoder-level domain invariance through an explicit paired objective rather than incidental exposure to degraded or limited real data, and verifies that this invariance -- not merely task accuracy -- survives cross-architecture distillation into hardware-deployable form.
The spatially and temporally co-registered $\Phi$-sat-2/Sentinel-2/simulated triplet dataset used in this thesis is what makes the first of these three possible; the paired-contrastive distillation objective developed in Chapter~\ref{chap:methodology} is what makes the second and third possible.
